\documentclass{article}

\usepackage[left=2cm,top=1.5cm,right=2cm,bottom=2cm]{geometry} 
\usepackage[utf8]{inputenc}
\usepackage[T1]{fontenc}
\usepackage[hidelinks]{hyperref}

% Resume Template by Cristian Palma Foster
\usepackage{graphicx} % manage images
\usepackage{wrapfig} % place images
\usepackage{multicol} % make columns
\usepackage{xcolor}
\usepackage{lmodern}
\pagenumbering{gobble} % erase page numbers

% Well spaced paragraphs
\setlength{\parskip}{-5pt}
\setlength\parindent{0pt}

% Well spaced lines in tables
\renewcommand{\arraystretch}{1.3}

% Contact Information
\newcommand{\contactinformation}[1]{
	\begin{center}
		{\LARGE \textbf{Cristian J. Palma Foster - #1}}

		\vspace{0.5cm}
		\href{mailto:cjp276@cornell.edu}{cjp276@cornell.edu},
		\href{http://www.linkedin.com/in/cristianpalmaf}{linkedin.com/in/cristianpalmaf},
		\href{https://cristianpalmafoster.github.io/}{cristianpalmafoster.github.io}
		% \href{http://www.github.com/cristianpalmafoster}{github.com/cristianpalmafoster}
		\vspace{0.2cm}
	\end{center}
}

\newenvironment{concept}[1]{
	\begin{large} \textcolor{gray}{\textbf{#1:}}\end{large}\\
	\vspace{-0.2cm}
	
}{}

\newcommand{\job}[2]{
	\normalsize \textbf{#1} \hfill \small \textcolor{gray}{\textit{#2}}\\
}
\newcommand{\data}[2]{
	\small #1 \hfill \small \textcolor{gray}{\textit{#2}}\\
}
\newcommand{\entry}[3]{
	\small #1 \textcolor{gray}{#2} \hfill \small \textcolor{gray}{\textit{#3}}\\
}
\newcommand{\desc}[1]{
	\small #1 
	\vspace{0.1cm}\\
}
\newcommand{\courses}[1]{
	\footnotesize \textcolor{gray}{#1}\\
}
\newcommand{\enumeration}[2]{
	\textcolor{gray}{\textbf{#1:}}
	{\small #2}
}


\begin{document}

\contactinformation

\enumeration{\large Investigación}{Teoría Algorítmica de Juegos, Aprendizaje en Juegos, Aprendizaje Online y Algoritmos de Aproximación.}

\begin{concept}{Educación}
	\institution{Universidad de Cornell}{Doctorado en Investigación de Operaciones. Profesora guía: Éva Tardos.}{2024 -- 2029}
	\annotation{Estadía de investigación con Anupam Gupta en la Universidad de Nueva York, 2026--27.}
	\institution{Universidad de Chile}{Magíster en Matemáticas Aplicadas. Profesor guía: José A. Soto.}{2022 -- 2023}
	% \annotation{Teoría de la complejidad, teoría de grafos, computación distribuida y algoritmos de aproximación.}
	\institution{Universidad de Chile y École CentraleSupélec}{Doble titulación en Ingeniería Matemática.}{2017 -- 2022}
	% \data{Licenciatura en Matemáticas e Ingeniería Informática, École CentraleSupélec.}{2019 -- 2021}
	% \annotation{Desarrollo web, propagación en grafos, redes y seguridad, algoritmos, procesamiento masivo de datos, aprendizaje automático y estadística.}
	% \data{Licenciatura en el Plan Común de Ingeniería, Universidad de Chile.}{2017 -- 2019}
\end{concept}

% Publicaciones y prepublicaciones
\begin{concept}{Artículos de Investigación}
	\job{\href{https://arxiv.org/abs/2502.08898}{Learning in Strategic Queuing Systems with Small Buffers}}{ACM EC 2026}
	\desc{Ariana Abel, Yoav Kolumbus, Jerónimo Martín Duque, Cristian Palma Foster y Éva Tardos.}
	\job{\href{https://arxiv.org/abs/2610.19179}{Learning in Strategic Queuing with Geometric Recharge: From Bandits to Systems}}{Prepublicación, 2026}
	\desc{Cristian Palma Foster y Éva Tardos.}
	\job{Maximizing Early Work via the Configuration LP}{Documento de trabajo}
	\desc{Cristian Palma Foster, Kevin Schewior y José Soto.}
	\job{Hidden Roles are as Expensive as Coloring}{Documento de trabajo}
	\desc{Cristian Palma Foster.}
	% \job{\href{https://repositorio.uchile.cl/handle/2250/200354}
	% {Algoritmos de balanceo de carga en problemas de agendamiento bajo una función cóncava de la carga}}
	% {Tesis de magíster, 2024}
	% \desc{Universidad de Chile. Profesor guía: José A. Soto. Escrita en español.}
\end{concept}

\begin{concept}{Charlas Seleccionadas} %y pósteres
	\talktopic{Sistemas de Colas Estratégicas}
	\data{Póster en el taller de NeurIPS sobre economía para el aprendizaje automático (EconML).}{Atlanta\enspace 2026}
	% \data{Revisión anual del proyecto MURI de la ONR, Arlington, Virginia}{Charla, 2026} %y póster
	\data{Conferencia de la ACM sobre Economía y Computación (EC).}{Roma\enspace 2026}
	% \data{Grupo de Algoritmos y Teoría de Juegos, Universidad McGill}{2026}
	% \data{Taller de Jóvenes Investigadores, Universidad de Cornell.}{Póster, 2025}
	% \data{Coloquio de doctorado de ORIE, Universidad de Cornell.}{Charla, 2025}
	\data{Seminario de Algoritmos, Teoría de Juegos, Combinatoria y Optimización (AGCO).}{Santiago\enspace 2025}
	\talktopic{Juegos con Roles Ocultos}
	\data{Póster en la Conferencia de la ACM sobre Economía y Computación (EC).}{Roma\enspace 2026}
	\talktopic{Balanceo de Carga con Objetivo Cóncavo}
	\data{Póster en la Jornada de Investigación de Operaciones de Nueva York, Universidad de Nueva York.}{Nueva York\enspace 2025}
	\data{Escuela de Verano en Matemática Discreta, Centro de Modelamiento Matemático.}{Valparaíso\enspace 2024} % XIX, Chile, y póster
	% \data{Simposio de Estudiantes de Posgrado de la Escuela de Ingeniería, Universidad de Chile.}{Charla y póster, 2023}
	\data{Sesión de Matemática Discreta, Reunión Anual de la Sociedad de Matemática de Chile (SOMACHI).}{Santiago\enspace 2023} % XCI 
	% \data{Seminario de Algoritmos, Teoría de Juegos, Combinatoria y Optimización (AGCO), Chile}{Charla, 2023}
	% \data{Seminario de Estudiantes de Magíster del Departamento de Ingeniería Matemática, Universidad de Chile.}{Chile, 2023}
\end{concept}

\begin{concept}{Experiencia Laboral}
	\job{Ingeniero de Investigación de Operaciones en el Centro de Modelamiento Matemático}{Santiago\enspace 2024}
	% \desc{Apliqué técnicas de programación lineal entera mixta a un problema de asignación de turnos de tripulación en Python.\\
	% Reduje de un día a minutos el tiempo de cálculo de los turnos de tripulación mediante programación lineal entera mixta en Python.}
	\desc{Reducción de un día a minutos el cálculo de los turnos de tripulación mediante un programa lineal mixto en Python.}
	\job{Asistente de Investigación en el Centro de Modelamiento Matemático}{Santiago\enspace 2023}%2022 -
	\desc{Desarrollo de algoritmos de aproximación para el balanceo de carga con objetivo cóncavo.}
	% \job{Profesor del Campeonato de Matemática CMAT}{Chile, 2022}
	% \desc{Preparé a estudiantes de enseñanza media para olimpiadas de matemáticas.}
	\job{Practicante de Ingeniería de Software en Bigblue}{París\enspace 2021}
	% \desc{Resolví problemas de logística de comercio electrónico y desarrollé servicios de backend en Golang con FoundationDB.} 
	\desc{Creación en Golang herramientas de envío y facturación, reduciendo en un 50~\% el tiempo de facturación.}
	% Habilité entregas de menor impacto ambiental al \textbf{integrar} un transportista ecológico al \textbf{sistema de gestión del transporte}.\\
	% Habilité el lanzamiento en España al integrar un transportista local.\\
	% Habilité la facturación en varios países al ampliar el \textbf{sistema de facturación} con configuraciones por país.\\
	% Reduje los errores de digitación en facturas mediante un \textbf{validador semántico} y un \textbf{detector de duplicados}.\\
	% Reduje el tiempo de facturación en un 50~\% al automatizar la \textbf{importación por lotes} de líneas de factura.}
	% \job{Práctica obrera de verano en Amazon}{Francia, 2020}
	% \desc{Asistí al encargado de bodega y me ocupé de la selección y el embalaje de pedidos.}
	% \job{Desarrollador del programa de intercambio de idiomas en el BdI de CentraleSupélec}{Francia, 2020}
	% \desc{Organicé encuentros de conversación en varios idiomas.\\
	% Emparejé a un 47~\% más de participantes mediante un \textbf{algoritmo de emparejamiento} en Python.}%Reduje el tiempo de cálculo a milisegundos
	% \job{Practicante de investigación en el Laboratorio de Procesamiento de Señales Sondra}{Francia, 2020}
	% \desc{Evalué un método de selección de orden para el \textbf{procesamiento de imágenes} y las \textbf{finanzas} en Matlab.\\
	% Aumenté en un 50~\% la rentabilidad de la inversión con una cartera de varianza mínima.}
	% \job{Académico del Campeonato de Matemática CMAT}{Chile, 2017 -- 2019}
	% \desc{Resolví problemas en YouTube y propuse problemas para las pruebas.}
\end{concept}

\begin{concept}{Honores}
	% \data{Reconocimiento como estudiante destacado de Ingeniería Matemática, Universidad de Chile.}{2021}
	\data{Beca de Excelencia Eiffel otorgada por el Ministerio para Europa y de Asuntos Exteriores.}{2019 -- 2021}
	\data{Primer lugar de la generación de Ingeniería, Universidad de Chile.}{2018}
	% \data{Primer lugar en la clasificación de la generación de Ingeniería, Universidad de Chile.}{2018}
	% \data{Reconocimiento como estudiante destacado del Plan Común de Ingeniería, Universidad de Chile.}{2017 -- 2018}
	\data{Beca de Excelencia Andrés Bello otorgada por la Universidad de Chile.}{2017 -- 2022}
	% \data{Octavo lugar en la admisión a Ingeniería, Universidad de Chile.}{2017}
	% \data{Puntaje nacional en las pruebas de admisión universitaria de Matemática y Ciencias, mención Física.}{2016}
	% \data{Reconocimiento al mejor rendimiento académico del Colegio Seminario Padre Alberto Hurtado.}{2016}
	\data{Medalla de oro en la Olimpiada Nacional de Matemática de la Sociedad de Matemática de Chile (SOMACHI).}{2016}
	% \data{Medalla de plata en la Olimpiada Nacional de Matemática de la SOMACHI, Chile.}{2015}
	% \data{Medalla de oro en la Olimpiada Regional de Física del Biobío de la SOCHIFI.}{2015}
	% \data{Medalla de bronce en la Olimpiada Nacional de Matemática de la SOMACHI, Chile.}{2014}
	% \data{Medalla de plata en la Olimpiada Nacional de Química de la SCHQ, Chile.}{2014}
	% \data{Medalla de bronce en la Olimpiada Nacional de Matemática de la SOMACHI, Chile.}{2013}
	% \data{Participación en la Olimpiada Nacional de Matemática de la SOMACHI, Chile.}{2012}
\end{concept}

\begin{concept}{Enseñanza Auxiliar} %Seleccionadas 
	% \data{\textbf{Teoría algorítmica de juegos} (2025) en la Universidad de Cornell}{}
	% \data{Herramientas prácticas para la \textbf{investigación de operaciones}, el \textbf{aprendizaje automático} y la \textbf{ciencia de datos} (2025) en la Universidad de Cornell}{}
	% \data{\textbf{Teoría de grafos} (2023) y \textbf{computabilidad y complejidad computacional} (2023) en la Universidad de Chile}{}
	% \data{\textbf{Programación lineal entera mixta} (2022, 2023) y \textbf{algoritmos combinatoriales} (2022) en la Universidad de Chile}{}
	% \data{Ecuaciones en derivadas parciales (2021), Teoría de la medida y probabilidad (2020) en École CentraleSupélec}{}
	% \data{Cálculo en varias variables (2019) y \textbf{algoritmos y estructuras de datos} (2018) en la Universidad de Chile}{}
	% \data{Cálculo diferencial e integral (2023, 2022) en la Universidad de Chile}{}
	% \data{Introducción al cálculo (2023, 2018) y Álgebra (2023, 2022) en la Universidad de Chile}{}
	\institution{Universidad de Cornell}{Teoría Algorítmica de Juegos, Herramientas Prácticas de IO, ML y Ciencia de Datos}{2025}
	\institution{Universidad de Chile}{Teoría de Grafos, Teoría de la Complejidad, Programación Lineal Mixta}{2023}
	\institution{École CentraleSupélec}{Ecuaciones en Derivadas Parciales, Teoría de la Medida y Probabilidad}{2021}
\end{concept}

\begin{concept}{Proyectos}
	\data{Balanceo de carga multirruta en redes en Python. Puesto número 15 en el desafío IMC de Huawei.}{2024}
	% \data{Desarrollé en Python protocolos de transmisión multirruta para el \textbf{balanceo de carga en redes}.}{2024}
	% \annotation{Obtuve el 15.º lugar en el desafío IMC de Huawei.}
	% \data{Desarrollé la \textbf{interfaz} en \textbf{JavaScript} de un planificador web de maratones de series con la \textbf{API} de TMDB.}{2021}%Binger
	% \data{Desarrollé la \textbf{interfaz} en \textbf{JavaScript} de un sitio de pronóstico del tiempo con la \textbf{API} de OpenWeather y \textbf{Node.js}.}{2021}%Aurora
	% \data{Optimicé la revaluación de seguros mediante \textbf{agrupamiento} y \textbf{descenso de gradiente} en Matlab con Generali.}{2021}%con Generali y CentraleSupélec
	\data{Algoritmo de intercambio de idiomas para CentraleSupélec emparejando en milisegundos a un 47~\% más de personas.}{2020}
	% \data{\textbf{Algoritmo de emparejamiento} para el intercambio de idiomas implementado en École CentraleSupélec.}{2020}
	% \annotation{Emparejé a un 47~\% más de participantes y reduje el tiempo de cálculo a milisegundos con una implementación en Python.}
	% \data{Control de ruido mediante \textbf{optimización paramétrica} de EDP en Python con ONERA.}{2020}%con ONERA y CentraleSupélec
	% \data{Desarrollé los \textbf{servicios de backend} de un sitio de recomendaciones de cine con Amazon DynamoDB.}{2020}%con Theodo y CentraleSupélec
	% \data{\textbf{Algoritmo de maximización de la propagación} en Python para una campaña de marketing viral con Artefact.}{2020}%con Artefact y CentraleSupélec
	% \data{Desarrollo integral de un prototipo de red social por proximidad geográfica para el hackatón de Facebook Chile.}{2019}%en JavaScript
	% \data{Desarrollo y montaje de un prototipo de robot guía turístico en Java con Lego Mindstorms.}{2018}%Ruteador
	% \data{Optimización de la multiplicación de matrices.}{2018}
	% \data{Compresión con códigos de Huffman.}{2018}
	% \data{Desarrollo y montaje de un prototipo de robot detector de piezas defectuosas en \textbf{Clang} con Arduino.}{2017}
	% \data{Réplica básica de un juego de Pokémon en \textbf{Python}.}{2017}
	% \data{Calculadora de expresiones y derivadas.}{2017}
\end{concept}

\enumeration{Programación}{Python, \LaTeX, Matlab, Julialang, Java, Clang, Golang, HTML y MySQL.}
% Programación: Python, Matlab, Julialang, Java, Clang y Golang.
% Redacción: MS Office, Jupyter y \LaTeX.
% Bases de datos: DynamoDB, FoundationDB y MySQL.
% Web: HTML, CSS, JavaScript, Node.js y React.
% Despliegue: Docker, Git, Google Cloud y Amazon Web Services.

\enumeration{Idiomas}{Inglés (Avanzado), Francés (Avanzado), Alemán (Elemental) y Español (Nativo).}

\enumeration{Tiempo Libre}{Vóleibol, Bodybuilding, Baile, Correr, Kung Fu y \textbf{Programación Competitiva}.}
% Kung fu (cinturón morado, dominio del bastón y del sable), carrera a pie (10 km de París en 52 min).
% Viajes
% Voluntariado: coordiné visitas a hogares de menores y residencias de adultos mayores. %Construcción de viviendas e implementación de programas de inclusión social.
\end{document}